\documentclass[12pt,a4paper]{article}
% ==============================================================================
% PREÁMBULO ESTANDARIZADO UNSCH - CIENCIAS DE LA SALUD (ENFERMERÍA)
% Compatible tanto con pdflatex como con xelatex (Unicode)
% ==============================================================================

\usepackage{iftex}

\ifPDFTeX
  \usepackage[utf8]{inputenc}
  \usepackage[T1]{fontenc}
  \usepackage{mathptmx} % Fuente Times 12pt requerida por la normativa
\else
  \usepackage{fontspec}
  \setmainfont{Times New Roman}
\fi

\usepackage[spanish,es-nodecimaldot,es-tabla]{babel}

% Márgenes normativos UNSCH:
% Superior: 3.0 cm | Izquierdo: 3.5 cm | Derecho: 2.5 cm | Inferior: 2.5 cm
\usepackage[top=3.0cm, left=3.5cm, right=2.5cm, bottom=2.5cm, a4paper]{geometry}

% Interlineado reglamentario (1.5 líneas)
\usepackage{setspace}
\onehalfspacing

% Sangría reglamentaria de párrafo
\setlength{\parindent}{1.27cm} % Sangría estilo APA (media pulgada)
\setlength{\parskip}{6pt}      % Separación ligera entre párrafos

% Paquetes matemáticos y símbolos
\usepackage{amsmath,amssymb,amsfonts}

% Tablas y matrices avanzadas
\usepackage{booktabs}
\usepackage{tabularx}
\usepackage{array}
\usepackage{multirow}
\usepackage{longtable}

% Inclusión de imágenes y gráficos
\usepackage{graphicx}
\usepackage{float}
\usepackage{caption}
\captionsetup{font=small, labelfont=bf, justification=centering}

% Encabezados y pie de página
\usepackage{fancyhdr}
\pagestyle{fancy}
\fancyhf{}
\fancyhead[R]{\small\thepage}
\fancyhead[L]{\small\nouppercase{\leftmark}}
\renewcommand{\headrulewidth}{0.5pt}
\renewcommand{\footrulewidth}{0pt}
\setlength{\headheight}{15pt}

% Estilo de páginas iniciales
\fancypagestyle{plain}{
    \fancyhf{}
    \fancyfoot[C]{\small\thepage}
    \renewcommand{\headrulewidth}{0pt}
}

% Formato de Títulos
\usepackage{titlesec}
\titleformat{\section}
    {\normalfont\large\bfseries}{\thesection.}{1em}{}
\titleformat{\subsection}
    {\normalfont\normalsize\bfseries}{\thesubsection.}{1em}{}
\titleformat{\subsubsection}
    {\normalfont\normalsize\itshape}{\thesubsubsection.}{1em}{}

% Control de listas numeradas y con viñetas
\usepackage{enumitem}
\setlist[itemize]{noitemsep, topsep=2pt, leftmargin=1.5cm}
\setlist[enumerate]{noitemsep, topsep=2pt, leftmargin=1.5cm}

% Compatibilidad con Pandoc
\providecommand{\tightlist}{%
  \setlength{\itemsep}{0pt}\setlength{\parskip}{0pt}}

% Colores y Enlaces (Modo Tinta Negra Oficial para Tesis y Protocolos UNSCH)
\usepackage{xcolor}

\usepackage{hyperref}
\hypersetup{
    colorlinks=true,
    linkcolor=black,
    citecolor=black,
    urlcolor=black,
    filecolor=black,
    pdfauthor={Gresly Lucero Pariona Palomino},
    pdftitle={Acceso a los servicios de Salud y Calidad de Atencion percibido por los usuarios},
    pdfsubject={Proyecto de Investigacion en Salud (EN 486) - UNSCH},
    pdfkeywords={Acceso a los servicios de salud, Calidad de atencion, Percepcion del usuario, Satisfaccion, Enfermeria, UNSCH, Ayacucho}
}

% Comillas en español
\usepackage{csquotes}

% Microtipografía
\usepackage{microtype}

% Compatibilidad con bloques de codigo de Pandoc
\ifx\Shaded\undefined
  \newenvironment{Shaded}{\begin{quote}\small\ttfamily}{\end{quote}}
\fi
\ifx\Highlighting\undefined
  \newenvironment{Highlighting}{}{}
\fi


\begin{document}
\section{CIENCIA Y CONOCIMIENTO}\label{ciencia-y-conocimiento}  \textbf{UNIVERSIDAD NACIONAL DE SAN CRISTÓBAL DE HUAMANGA}\\ \textbf{FACULTAD DE CIENCIAS DE LA SALUD}\\ \textbf{ESCUELA PROFESIONAL DE ENFERMERÍA}  \begin{itemize} \tightlist \item   \textbf{Asignatura:} Proyecto de Investigación en Salud (EN 486)\\ \item   \textbf{Docente:} Dr.~Manglio Aguirre Andrade    \begin{itemize}   \tightlist   \item     \emph{Docente en la Categoría Principal de la UNSCH}\\   \item     \emph{Experto en Salud Pública}\\   \item     \textbf{Correo:} manglio.aguirre@unsch.edu.pe\\   \item     \textbf{Teléfono:} 966102220   \end{itemize} \end{itemize}  \begin{center}\rule{0.5\linewidth}{0.5pt}\end{center}  \subsection{1. INTRODUCCIÓN}\label{introducciuxf3n}  \subsubsection{¿A qué nos enfrentamos?}\label{a-quuxe9-nos-enfrentamos}  \begin{itemize} \tightlist \item   Pandemia. \item   Crisis económica mundial. \item   Recursos limitados o insuficientes. \item   Cambio climático. \item   Presión social. \item   Conflictos armados. \item   Violencia social. \item   Sobredemanda de servicios de atención. \end{itemize}  \subsubsection{Tendencias y Desafíos en Salud}\label{tendencias-y-desafuxedos-en-salud}  \begin{itemize} \tightlist \item   Tendencias crecientes en demanda de servicios médicos incluyendo toda   clase de tecnologías en salud. \item   Cambios sociodemográficos: 9,4\% de la población mundial   \textgreater{} 65 años en 2020. \item   Prevalencia de enfermedades No Transmisibles y Transmisibles. \item   Presencia cada vez más relevante de la comunidad exigiendo sus   derechos. \item   Avances de la genética y mapeo del genoma humano. \item   Desarrollo de la tecnología y la informática. \item   ¿Control de los mercados y fin del libre comercio? \end{itemize}  \begin{center}\rule{0.5\linewidth}{0.5pt}\end{center}  \subsection{2. CIENCIA: GENERALIDADES}\label{ciencia-generalidades}  \subsubsection{Definición de Ciencia}\label{definiciuxf3n-de-ciencia}  En general, la ciencia es un conocimiento \textbf{racional, sistemático, exacto, verificable y por consiguiente falible}.\\ Por medio de la ciencia, el hombre ha alcanzado una reconstrucción conceptual del mundo cada vez más amplia, profunda y exacta.  \begin{quote} \textbf{EPISTEMOLOGÍA:}\\ Disciplina que se ocupa de estudiar los métodos que se emplean para alcanzar el conocimiento científico y las formas de validar dicho conocimiento. \end{quote}  \subsubsection{Perspectivas Teóricas}\label{perspectivas-teuxf3ricas}  \paragraph{Edgar Morin (Sociólogo Francés)}\label{edgar-morin-sociuxf3logo-francuxe9s}  \begin{itemize} \tightlist \item   La ciencia nos lleva a comprender la complejidad de la realidad. \item   Para analizar los problemas reales se necesitan herramientas   conceptuales. \item   Las herramientas teórico-metodológicas nos permiten construir el nuevo   conocimiento. \end{itemize}  \paragraph{Roszak}\label{roszak}  \begin{quote} \emph{``Los que queremos alcanzar un mejor conocimiento de la realidad que nos rodea nos calificamos de científicos.''} \end{quote}  \paragraph{Descubrimientos por Serendipia (Casualidad)}\label{descubrimientos-por-serendipia-casualidad}  Muchos de los descubrimientos que cambiaron nuestras vidas fueron hallados de casualidad: * \textbf{Alexander Fleming:} Descubrió la penicilina y salvó con ello más vidas. Al irse de vacaciones dejó su laboratorio sucio; al regresar, observó en el microscopio lo ocurrido en sus mezclas y descubrió el hongo \emph{Penicillium} que mataba a las bacterias. * \textbf{Percy Spencer:} Mucho menos conocido, pero probablemente tengas su invento en tu casa. Trabajando con un magnetrón se dio cuenta de que el chocolate en su bolsillo se había derretido. A partir de ello inventó el horno microondas \emph{(Pérez, L., Pérez, R. y Seca, M.V., 2020)}.  \subsubsection{La Ciencia como Sistema (Mario Bunge)}\label{la-ciencia-como-sistema-mario-bunge}  Según \textbf{Bunge, M. (1980)}, toda ciencia forma parte de un sistema de ciencias, en el sentido de que cada una de estas ciencias tiene alguna ciencia vecina con la que se solapa aunque sea parcialmente.  Podemos relacionar la ciencia como una \textbf{decatupla}, es decir, como un conjunto de diez elementos contextuales que determinan el desarrollo del conocimiento:  \[\text{CP} = \langle C, S, D, G, F, B, P, A, O, M \rangle\]  \begin{itemize} \tightlist \item   \textbf{(C) Comunidad:} Comunidad científica. \item   \textbf{(S) Sociedad:} Contexto social en el que se desarrolla. \item   \textbf{(D) Dominio:} Universo de discurso / objetos de estudio. \item   \textbf{(G) Supuestos filosóficos:} Supuestos ontológicos y   gnoseológicos de la ciencia. \item   \textbf{(F) Fondo Formal:} Herramientas lógicas y matemáticas. \item   \textbf{(B) Fondo Específico:} Teorías y conocimientos de campos   afines. \item   \textbf{(P) Problemática:} Problemas que se propone investigar. \item   \textbf{(A) Fondo Acumulado:} Fondo de conocimientos acumulado   previamente. \item   \textbf{(O) Objetivos:} Los fines y metas de la ciencia. \item   \textbf{(M) Metódica:} Colección de métodos y técnicas empleados. \end{itemize}  \begin{center}\rule{0.5\linewidth}{0.5pt}\end{center}  \subsection{2. CIENCIA: CARACTERÍSTICAS (Bunge, M.)}\label{ciencia-caracteruxedsticas-bunge-m.}  \begin{enumerate} \def\labelenumi{\arabic{enumi}.} \tightlist \item   \textbf{El conocimiento científico es fáctico:} Parte de los hechos,   los respeta hasta cierto punto y siempre vuelve a ellos. \item   \textbf{El conocimiento científico trasciende los hechos:} Descarta   hechos, produce nuevos hechos y los explica. \item   \textbf{La ciencia es analítica:} Aborda problemas circunscritos, uno   a uno, y trata de descomponerlo todo en elementos. \item   \textbf{La investigación científica es especializada:} Consecuencia   del enfoque analítico de los problemas. \item   \textbf{El conocimiento científico es claro y preciso:} Sus problemas   son distintos, sus resultados son claros. \item   \textbf{El conocimiento científico es comunicable:} No es inefable ni   privado, sino expresable y público. \item   \textbf{El conocimiento científico es verificable:} Debe superar el   examen de la experiencia (testabilidad). \item   \textbf{La investigación científica es metódica:} No es errática sino   planeada. \item   \textbf{El conocimiento científico es sistemático:} Es un sistema de   ideas conectadas lógicamente entre sí. \item   \textbf{El conocimiento científico es general:} Ubica los hechos   singulares en pautas generales. \item   \textbf{El conocimiento científico es legal:} Busca leyes de la   naturaleza y de la sociedad y las aplica. \item   \textbf{La ciencia es explicativa:} Intenta explicar los hechos en   términos de leyes, y las leyes en términos de principios. \item   \textbf{El conocimiento científico es predictivo:} Trasciende la masa   de los hechos de experiencia, imaginando cómo pudo haber sido el   pasado y cómo podrá ser el futuro. \item   \textbf{La ciencia es abierta:} No reconoce barreras a priori que   limiten el conocimiento; es falible y capaz de autocorregirse. \item   \textbf{La ciencia es útil:} Porque busca la verdad, la ciencia es   eficaz en la provisión de herramientas para el bien o para el mal. \end{enumerate}  \begin{center}\rule{0.5\linewidth}{0.5pt}\end{center}  \subsection{2. CIENCIA: CLASIFICACIÓN}\label{ciencia-clasificaciuxf3n}  \subsubsection{Clasificación según Erik Wolf}\label{clasificaciuxf3n-seguxfan-erik-wolf}  \begin{itemize} \tightlist \item   \textbf{Según su importancia en orden a la fundamentación del   conocimiento científico:}    \begin{itemize}   \tightlist   \item     Ciencias fundamentales generales.   \item     Ciencias fundamentales especiales.   \end{itemize} \item   \textbf{Según su objeto o método de conocimiento:}    \begin{itemize}   \tightlist   \item     Ciencias de la naturaleza.   \item     Ciencias del espíritu o de la cultura.   \end{itemize} \item   \textbf{Según su finalidad:}    \begin{itemize}   \tightlist   \item     Ciencias teóricas.   \item     Ciencias prácticas.   \end{itemize} \end{itemize}  \begin{center}\rule{0.5\linewidth}{0.5pt}\end{center}  \subsubsection{Clasificación según Mario Bunge}\label{clasificaciuxf3n-seguxfan-mario-bunge}  \begin{verbatim}                            CIENCIA                               │              ┌────────────────┴────────────────┐              ▼                                 ▼           FORMAL                            FACTUAL       ┌──────┴──────┐                   ┌──────┴──────┐       ▼             ▼                   ▼             ▼    Lógica      Matemáticas           NATURAL      ANTRÓPICAS (Sociales)                                   ┌─────┼─────┬─────┐   ┌─────┼─────┬─────┐                                   ▼     ▼     ▼     ▼   ▼     ▼     ▼     ▼                                Física Química Bio. Psic. Soc. Econ. C.Pol. Hist. \end{verbatim}  \paragraph{a. Ciencias Formales}\label{a.-ciencias-formales}  No se ocupan de hechos fácticos, sino de abstracciones creadas en nuestro pensamiento, como la lógica y la matemática. Tratan de entes ideales, es decir, de entes que sólo existen en la mente humana. Por eso, a los lógicos y matemáticos no se les dan los objetos de estudio, sino que ellos los construyen.  \paragraph{b. Ciencias Fácticas}\label{b.-ciencias-fuxe1cticas}  Se ocupan de hechos que ocurren fuera o dentro de nuestros cerebros, los cuales son cosas concretas o materiales. Por lo tanto, estos hechos tienen propiedades físicas, biológicas y sociales.  \begin{center}\rule{0.5\linewidth}{0.5pt}\end{center}  \subsubsection{Diferencia entre Ciencias Formales y Fácticas}\label{diferencia-entre-ciencias-formales-y-fuxe1cticas}  {\def\LTcaptype{none} % do not increment counter \begin{longtable}[]{@{}   >{\raggedright\arraybackslash}p{(\linewidth - 4\tabcolsep) * \real{0.3333}}   >{\raggedright\arraybackslash}p{(\linewidth - 4\tabcolsep) * \real{0.3333}}   >{\raggedright\arraybackslash}p{(\linewidth - 4\tabcolsep) * \real{0.3333}}@{}} \toprule\noalign{} \begin{minipage}[b]{\linewidth}\raggedright Conceptos Ordenadores \end{minipage} & \begin{minipage}[b]{\linewidth}\raggedright Ciencias Formales \end{minipage} & \begin{minipage}[b]{\linewidth}\raggedright Ciencias Fácticas \end{minipage} \\ \midrule\noalign{} \endhead \bottomrule\noalign{} \endlastfoot \textbf{Naturaleza del objeto} & Objetos formales (ideales, abstractos) & Objetos factuales (materiales, concretos) \\ \textbf{Método de comprobación de proposiciones} & Por demostración o deducción & Por verificación (observación / experimentación) \\ \textbf{Criterio de verdad} & Coherencia lógica & Correspondencia objetiva con los hechos \\ \textbf{Carácter de los enunciados} & Lógicamente necesarios & Lógicamente necesarios y verificables empíricamente \\ \end{longtable} }  \begin{center}\rule{0.5\linewidth}{0.5pt}\end{center}  \subsubsection{Clasificación por Finalidad y Aplicación}\label{clasificaciuxf3n-por-finalidad-y-aplicaciuxf3n}  \begin{itemize} \tightlist \item   \textbf{Ciencias Básicas o Teóricas:} El objeto de estudio es   desarrollar nuevas teorías, fortalecer o refutar las existentes; trata   de comprender los elementos esenciales y la relación entre ellos   dentro de un contexto determinado. \item   \textbf{Ciencias Aplicadas o Prácticas:} Las verdades que buscan son   de posible utilización práctica o valor aplicativo. \emph{Ejemplo:} La   investigación de productos naturales para determinar si pueden   utilizarse en la industria farmacéutica. \item   \textbf{Técnica:} El técnico, a diferencia del científico, asigna   valores a todas las cosas. La técnica se refiere al hacer eficaz, es   decir, a reglas que permitan alcanzar de modo correcto, preciso y   satisfactorio ciertos objetivos prácticos \emph{(Santiesteban Naranjo,   E., 2014)}. \end{itemize}  \begin{center}\rule{0.5\linewidth}{0.5pt}\end{center}  \subsubsection{Funciones Principales de la Ciencia Contemporánea}\label{funciones-principales-de-la-ciencia-contemporuxe1nea}  \emph{(Santiesteban Naranjo, E., 2014)} 1. \textbf{Cognoscitiva:} Búsqueda del saber y comprensión de la realidad. 2. \textbf{Práctica:} Aplicación transformadora del conocimiento en beneficio de la sociedad. 3. \textbf{Formativa:} Desarrollo del pensamiento crítico e investigativo. 4. \textbf{Educativa:} Transmisión y difusión cultural del conocimiento.  \begin{center}\rule{0.5\linewidth}{0.5pt}\end{center}  \subsection{3. EL CONOCIMIENTO}\label{el-conocimiento}  \subsubsection{Definición}\label{definiciuxf3n}  Es la aprehensión o captación de la imagen de un objeto. Está constituida por un conjunto de cualidades o propiedades de ese objeto. La ``aprehensión'' o captura de las cualidades es mental; no es una captura física.  \begin{quote} \textbf{Karl Popper:}\\ El conocimiento científico es una aproximación crítica a la realidad, apoyándose en el estudio del método científico que, fundamentalmente, trata de percibir y explicar desde lo esencial hasta lo más complejo: el porqué de las cosas y su devenir. \end{quote}  \subsubsection{Relación Sujeto - Objeto en el Conocimiento Científico}\label{relaciuxf3n-sujeto---objeto-en-el-conocimiento-cientuxedfico}  \begin{verbatim}               ┌───────────────────────────────┐               │   Obtener información,        │               │   conseguir datos             │               └──────────────┬────────────────┘                              │     ┌──────────┐      INTERRELACIÓN      ┌──────────┐     │  SUJETO  │ ◄─────────────────────► │  OBJETO  │     └────┬─────┘                         └─────┬────┘          │                                     │          └──────────────────┬──────────────────┘                             ▼                      [ CONOCIMIENTO ] \end{verbatim}  \begin{center}\rule{0.5\linewidth}{0.5pt}\end{center}  \subsubsection{Tipos de Conocimiento}\label{tipos-de-conocimiento}  \begin{verbatim}                                CONOCIMIENTO                                     │            ┌────────────────────────┴────────────────────────┐            ▼                                                 ▼ Conocimiento No Científico                         Conocimiento Científico            │                                                 │ ┌──────────┴──────────┐                           ┌──────────┴──────────┐ ▼                     ▼                           ▼                     ▼ Empírico espontáneo   Mítico - Religioso          Empírico              Teórico (común, cotidiano, ordinario) \end{verbatim}  \paragraph{1. Conocimiento Empírico}\label{conocimiento-empuxedrico}  \begin{itemize} \tightlist \item   Su contenido surge de la experiencia con cierta elaboración racional. \item   El reflejo del objeto está limitado a la contemplación sensorial y los   datos que se obtienen; es la base práctica para la elaboración del   conocimiento teórico. \item   Es un modo de conocer adquirido en el proceso de socialización, desde   que somos niños y vamos aprendiendo cosas, y que hace que cada grupo   humano despliegue modos particulares de explorar la realidad   \emph{(Pérez, L., Pérez, R. y Seca, M.V., 2020)}. \end{itemize}  \paragraph{2. Conocimiento Mítico - Religioso}\label{conocimiento-muxedtico---religioso}  \begin{itemize} \tightlist \item   Este saber parte de una verdad externa a la realidad observada, de la   que se deducen todos los fenómenos constitutivos de lo real. \item   Tiene una pretensión totalizante, en tanto subsume la explicación de   todos los fenómenos en la verdad de sus principios. \item   Tiene carácter dogmático (no se discute).\\   \emph{Ejemplo:} ``El ser humano fue creado por Dios, como la tierra y   el cielo y la naturaleza. Por lo tanto, no importa, por ejemplo, que   los antropólogos planteen la teoría de la evolución'' \emph{(Pérez,   L., Pérez, R. y Seca, M.V., 2020)}. \end{itemize}  \paragraph{3. Conocimiento Teórico (Científico)}\label{conocimiento-teuxf3rico-cientuxedfico}  \begin{itemize} \tightlist \item   Es el nivel de conocimiento que se obtiene sobre el objeto, a partir   de la experiencia, mediante el pensamiento abstracto \emph{(Hernández   León, R.A., 2011)}. \item   Tiene un carácter convencional, ya que los métodos que se siguen para   producir saberes y el lenguaje que se utiliza para comunicarlos se   basan en acuerdos generados por la comunidad científica. \item   Es producto de la aplicación del método científico. No es dogmático:   cambia constantemente; incluso los saberes que parecen más firmes   (como la estructura del átomo o la tabla periódica) evolucionan con el   tiempo \emph{(Pérez, L., Pérez, R. y Seca, M.V., 2020)}. \end{itemize}  \begin{center}\rule{0.5\linewidth}{0.5pt}\end{center}  \subsubsection{Características del Conocimiento Científico}\label{caracteruxedsticas-del-conocimiento-cientuxedfico}  \begin{itemize} \tightlist \item   \textbf{Objetivo:} Se ajusta a los hechos de la realidad. \item   \textbf{Racional:} Utiliza la razón, conceptos, juicios y   razonamientos lógicos. \item   \textbf{Fundamentación:} Se apoya en demostraciones, pruebas objetivas   o documentos. \item   \textbf{Necesidad:} Un conocimiento es necesario cuando es válido para   todas las épocas y todos los lugares, cuando no varía de un tiempo a   otro, de un lugar a otro. \item   \textbf{Universalidad:} Es aquel que es válido para todos los seres   humanos. \item   \textbf{Sistemático:} Constituye un sistema orgánico y coherente de   ideas interrelacionadas. \item   \textbf{Crítico:} Se somete a constante examen, juicio y evaluación de   sus fundamentos. \item   \textbf{Comunicable:} Se expresa mediante un lenguaje claro y   accesible a la comunidad científica. \end{itemize}  \begin{center}\rule{0.5\linewidth}{0.5pt}\end{center}  \subsubsection{Elementos del Conocimiento}\label{elementos-del-conocimiento}  El acto de conocer articula cuatro componentes esenciales: 1. \textbf{Sujeto cognoscente:} Quien realiza el acto de conocer. 2. \textbf{Objeto cognoscible:} Aquello que es percibido o aprehendido. 3. \textbf{La operación misma de conocer:} El proceso psíquico y mental de captación del objeto. 4. \textbf{El resultado obtenido:} El \textbf{conocimiento mismo} o la representación conceptual del \textbf{objeto conocido}.  \begin{center}\rule{0.5\linewidth}{0.5pt}\end{center}  \subsubsection{Fases del Conocimiento}\label{fases-del-conocimiento}  \begin{verbatim}                     FASES DEL CONOCIMIENTO                               │              ┌────────────────┴────────────────┐              ▼                                 ▼       FASE SENSORIAL                     FASE RACIONAL              │                                 │      1. Sensación                      1. Concepto      2. Percepción                     2. Juicio      3. Representación                 3. Razonamiento \end{verbatim}  \begin{itemize} \tightlist \item   \textbf{Fase Sensorial:}    \begin{itemize}   \tightlist   \item     \textbf{Sensación:} Reflejo de cualidades aisladas del objeto a     través de los sentidos.   \item     \textbf{Percepción:} Reflejo integral del objeto en la conciencia     sensible.   \item     \textbf{Representación:} Imagen mental evocada en ausencia del     objeto concreto.   \end{itemize} \item   \textbf{Fase Racional:}    \begin{itemize}   \tightlist   \item     \textbf{Concepto:} Abstracción y generalización de los rasgos     esenciales del objeto.   \item     \textbf{Juicio:} Afirmación o negación de una relación entre     conceptos.   \item     \textbf{Razonamiento:} Deducción o inferencia de nuevos     conocimientos a partir de juicios previos.   \end{itemize} \end{itemize}  \begin{center}\rule{0.5\linewidth}{0.5pt}\end{center}  \subsubsection{El Conocimiento Científico frente a la Realidad}\label{el-conocimiento-cientuxedfico-frente-a-la-realidad}  \begin{verbatim}     [ REALIDAD ] ──► OBSERVA ──► DESCUBRE ──► EXPLICA ──► PREDICE ──► [ CONOCIMIENTO SISTEMÁTICO DE LA REALIDAD ] \end{verbatim}  A través de la observación metódica de la realidad, el científico descubre regularidades, explica sus causas y mecanismos, y formula predicciones sobre su comportamiento futuro, consolidando un conocimiento sistemático.  \begin{center}\rule{0.5\linewidth}{0.5pt}\end{center}  \subsection{4. PROBLEMAS DEL CONOCIMIENTO}\label{problemas-del-conocimiento}  La epistemología identifica cinco problemas principales en la teoría del conocimiento:  \begin{enumerate} \def\labelenumi{\arabic{enumi}.} \tightlist \item   \textbf{La posibilidad del conocimiento humano:}\\   ¿Puede realmente el sujeto aprehender el objeto? \item   \textbf{El origen del conocimiento:}\\   ¿Es la razón o la experiencia la fuente principal del conocimiento   humano? \item   \textbf{La esencia del conocimiento humano:}\\   ¿Es el objeto quien determina al sujeto o es el sujeto quien determina   al objeto? \item   \textbf{Las formas del conocimiento humano:}\\   ¿El conocimiento es exclusivamente racional o puede ser también   intuitivo? \item   \textbf{El criterio de verdad:}\\   ¿Cómo sabemos que nuestro conocimiento es verdadero y qué nos   garantiza su certeza? \end{enumerate}  \begin{center}\rule{0.5\linewidth}{0.5pt}\end{center}  \subsection{4. APLICACIÓN DEL CONOCIMIENTO}\label{aplicaciuxf3n-del-conocimiento}  El desarrollo científico y el conocimiento generado en el ámbito de la salud se orientan a transformar las prácticas asistenciales, fortalecer la salud pública y brindar soluciones concretas a las necesidades de la comunidad.
\end{document}
